% GENERATED FILE. Edit canonical lessons, then run npm run generate:books.
% canonical-source-hash: fnv1a64:2d7eeecea49d9d21
% canonical-lessons: MR-C277-jain, MR-C277-janar-ahe, MR-R277-what-will-happen

\chapter{Will, Going To, Will Not}
\label{ch:mr-will-going-to-will-not}
\hlchaptermodality{277}

\begin{chapteropening}
\textbf{By the end of this chapter:} \emph{I can say what will happen (jāīn, paḍel), what I am going to do (jāṇār āhe) and what I won't do (jāṇār nāhī).}

Complete MR-R277-what-will-happen: say what will happen, what is planned and what won't happen.
\end{chapteropening}

\section[\texorpdfstring{\mr{मी} \mr{जाईन} (mī jāīn)}{मी जाईन (mī jāīn)}]{\mr{मी} \mr{जाईन} (mī jāīn) — the future, person by person}

\label{lesson:MR-C277-jain}



\begin{quote}
Say \textquotedblleft{}see you tomorrow\textquotedblright{}: \emph{udyā bheṭū}. Then \textquotedblleft{}I went\textquotedblright{}: \emph{mī gelo}.

\textbf{\mr{भेटू}}, \textquotedblleft{}we'll meet\textquotedblright{}, has been a future all along. Here are the others.
\end{quote}

\subsection*{You'll want to know: the future, person by person}
\begin{quote}
\textbf{\mr{उद्या} \mr{मी} \mr{पुण्याला} \mr{जाईन}.}
\end{quote}

\begin{quote}
\emph{udyā mī puṇyālā jāīn.} — \textquotedblleft{}Tomorrow I will go to Pune.\textquotedblright{}
\end{quote}

\begin{itemize}
\raggedright
  \item \textbf{\mr{मी}} \textbf{\mr{जाईन}} (\emph{jāīn}) — I will go
  \item \textbf{\mr{तू}} \textbf{\mr{जाशील}} (\emph{jāśīl}) — you will go, to a friend
  \item \textbf{\mr{तो}}, \textbf{\mr{ती}}, \textbf{\mr{ते}} \textbf{\mr{जाईल}} (\emph{jāīl}) — he, she, it will go
  \item \textbf{\mr{आम्ही}} \textbf{\mr{जाऊ}} (\emph{jāū}) — we will go: the \textbf{-\mr{ऊ}} of \textbf{\mr{भेटू}}
  \item \textbf{\mr{तुम्ही}} \textbf{\mr{जाल}} (\emph{jāl}) — you will go, with respect
  \item \textbf{\mr{ते}} \textbf{\mr{जातील}} (\emph{jātīl}) — they will go
\end{itemize}

There is no gender here: a man and a woman both say \textbf{\mr{मी} \mr{जाईन}}.

\subsection*{You'll want to know: more verbs in the future}
The endings stay; the stem in front changes:

\begin{itemize}
\raggedright
  \item \textbf{\mr{येणे}} $\to$ \textbf{\mr{मी} \mr{येईन}} (\emph{mī yeīn}), \textbf{\mr{तो} \mr{येईल}} (\emph{to yeīl}), \textbf{\mr{तुम्ही} \mr{याल}} (\emph{tumhī yāl})
  \item \textbf{\mr{करणे}} $\to$ \textbf{\mr{मी} \mr{करीन}} (\emph{mī karīn}), \textbf{\mr{तो} \mr{करेल}} (\emph{to karel}), \textbf{\mr{आम्ही} \mr{करू}} (\emph{āmhī karū})
  \item \textbf{\mr{वाचणे}} $\to$ \textbf{\mr{मी} \mr{वाचीन}} (\emph{mī vācīn}), \textbf{\mr{ती} \mr{वाचेल}} (\emph{tī vācel})
  \item \textbf{\mr{पडणे}} $\to$ \textbf{\mr{पाऊस} \mr{पडेल}} (\emph{pāūs paḍel}), \textquotedblleft{}it will rain\textquotedblright{}
  \item \textbf{\mr{असणे}} $\to$ \textbf{\mr{मी} \mr{घरी} \mr{असेन}} (\emph{mī gharī asen}), \textquotedblleft{}I will be at home\textquotedblright{}; \textbf{\mr{तो} \mr{असेल}} (\emph{to asel})
  \item \textbf{\mr{होणे}} $\to$ \textbf{\mr{काय} \mr{होईल}?} (\emph{kāy hoīl?}), \textquotedblleft{}What will happen?\textquotedblright{}
\end{itemize}

Many people say \textbf{\mr{करेन}} (\emph{karen}) and \textbf{\mr{वाचेन}} (\emph{vācen}) for \textquotedblleft{}I will\textquotedblright{}; both are heard.

\subsection*{Your turn}
Say these aloud:

\begin{itemize}
\raggedright
  \item \textquotedblleft{}Tomorrow I will go to Pune\textquotedblright{} — \emph{udyā mī puṇyālā jāīn}
  \item \textquotedblleft{}It will rain in the evening\textquotedblright{} — \emph{sandhyākāḷī pāūs paḍel}
  \item \textquotedblleft{}We will meet tomorrow\textquotedblright{} — \emph{āmhī udyā bheṭū}
  \item \textquotedblleft{}What will happen?\textquotedblright{} — \emph{kāy hoīl?}
\end{itemize}

\begin{tcolorbox}[breakable,colback=teal!4,colframe=teal!35!black,title={Before you move on}]
\textquotedblleft{}I will go\textquotedblright{}? (\textbf{\mr{जाईन}}.) \textquotedblleft{}They will go\textquotedblright{}? (\textbf{\mr{जातील}}.) \textquotedblleft{}We will go\textquotedblright{}? (\textbf{\mr{जाऊ}}.) \textquotedblleft{}It will rain\textquotedblright{}? (\textbf{\mr{पाऊस} \mr{पडेल}}.) Does a woman say it differently? (No: the future has no gender.)
\end{tcolorbox}
\section[\texorpdfstring{\mr{जाणार} \mr{आहे} (jāṇār āhe)}{जाणार आहे (jāṇār āhe)}]{\mr{जाणार} \mr{आहे} (jāṇār āhe) — a plan, and won't}

\label{lesson:MR-C277-janar-ahe}



\begin{quote}
Say \textquotedblleft{}I will go\textquotedblright{}. (\emph{mī jāīn}.) Then \textquotedblleft{}next week\textquotedblright{}. (\emph{puḍhcyā āṭhavḍyāt}.)
\end{quote}

\subsection*{You'll want to know: -\mr{णार} \mr{आहे}, 'going to'}
\begin{quote}
\textbf{\mr{पुढच्या} \mr{आठवड्यात} \mr{मी} \mr{गावाला} \mr{जाणार} \mr{आहे}.}
\end{quote}

\begin{quote}
\emph{puḍhcyā āṭhavḍyāt mī gāvālā jāṇār āhe.} — \textquotedblleft{}Next week I am going to go to the village.\textquotedblright{}
\end{quote}

Take \textbf{-\mr{णे}} off the dictionary form and put \textbf{-\mr{णार}} on: \textbf{\mr{जाणे}} $\to$ \textbf{\mr{जाणार}} (\emph{jāṇār}), \textbf{\mr{करणे}} $\to$ \textbf{\mr{करणार}} (\emph{karṇār}). Then the present of \textquotedblleft{}to be\textquotedblright{} carries the person: \textbf{\mr{मी} \mr{जाणार} \mr{आहे}} (\emph{mī jāṇār āhe}), \textbf{\mr{तुम्ही} \mr{जाणार} \mr{आहात}} (\emph{tumhī jāṇār āhāt}), and, with \textbf{\mr{आहोत}} (\emph{āhot}, \textquotedblleft{}we are\textquotedblright{}), \textbf{\mr{आम्ही} \mr{जाणार} \mr{आहोत}} (\emph{āmhī jāṇār āhot}). It says a plan already made, where \textbf{\mr{जाईन}} can be a promise made on the spot.

\subsection*{You'll want to know: -\mr{णार} \mr{नाही}, 'won't'}
\begin{quote}
\textbf{\mr{मी} \mr{उद्या} \mr{येणार} \mr{नाही}.}
\end{quote}

\begin{quote}
\emph{mī udyā yeṇār nāhī.} — \textquotedblleft{}I won't come tomorrow.\textquotedblright{}
\end{quote}

Put \textbf{\mr{नाही}} in place of \textbf{\mr{आहे}}. This is the everyday \textquotedblleft{}won't\textquotedblright{} for every person: \textbf{\mr{तो} \mr{जाणार} \mr{नाही}} (\emph{to jāṇār nāhī}), \textbf{\mr{आम्ही} \mr{करणार} \mr{नाही}} (\emph{āmhī karṇār nāhī}).

\subsection*{Your turn}
Say these aloud:

\begin{itemize}
\raggedright
  \item \textquotedblleft{}Next week I am going to the village\textquotedblright{} — \emph{puḍhcyā āṭhavḍyāt mī gāvālā jāṇār āhe}
  \item \textquotedblleft{}I won't come tomorrow\textquotedblright{} — \emph{mī udyā yeṇār nāhī}
  \item \textquotedblleft{}What are you going to do?\textquotedblright{} — \emph{tumhī kāy karṇār āhāt?}
\end{itemize}

\begin{tcolorbox}[breakable,colback=teal!4,colframe=teal!35!black,title={Before you move on}]
\textquotedblleft{}I am going to go\textquotedblright{}? (\textbf{\mr{मी} \mr{जाणार} \mr{आहे}}.) \textquotedblleft{}I won't go\textquotedblright{}? (\textbf{\mr{मी} \mr{जाणार} \mr{नाही}}.) \textquotedblleft{}What are you going to do?\textquotedblright{} (\textbf{\mr{तुम्ही} \mr{काय} \mr{करणार} \mr{आहात}?})
\end{tcolorbox}
\section[(review)]{Review — what will happen}

\label{lesson:MR-R277-what-will-happen}



\begin{quote}
Say \textquotedblleft{}Yesterday I went to Pune\textquotedblright{} and \textquotedblleft{}Tomorrow I will go to Pune\textquotedblright{}. (\emph{kāl mī puṇyālā gelo}, \emph{udyā mī puṇyālā jāīn}.) Then \textquotedblleft{}he did the work\textquotedblright{}. (\emph{tyāne kām kelaṁ}.)
\end{quote}

\subsection*{Your turn}
Say these aloud:

\begin{itemize}
\raggedright
  \item Yesterday it rained. Tomorrow it will rain too.
  \item Yesterday I was not at home. Tomorrow I will be at home.
  \item Next week she is going to go to the village.
  \item I won't come tomorrow, because I am going to read a book.
  \item What will happen? What happened?
\end{itemize}

\begin{tcolorbox}[breakable,colback=teal!4,colframe=teal!35!black,title={Before you move on}]
\textquotedblleft{}I will\textquotedblright{}? (\textbf{-\mr{ईन}}: \textbf{\mr{जाईन}}, \textbf{\mr{येईन}}.) \textquotedblleft{}It will\textquotedblright{}? (\textbf{-\mr{एल}} or \textbf{-\mr{ईल}}: \textbf{\mr{पडेल}}, \textbf{\mr{जाईल}}.) A plan? (\textbf{-\mr{णार} \mr{आहे}}.) \textquotedblleft{}Won't\textquotedblright{}? (\textbf{-\mr{णार} \mr{नाही}}.)
\end{tcolorbox}
